\documentclass[12pt, a4paper]{article}

\usepackage[utf8]{inputenc}
\usepackage{amsfonts, amsmath, amssymb, amsthm}
\usepackage{geometry}
\usepackage[hidelinks]{hyperref}
\usepackage[none]{hyphenat}
\usepackage{setspace}
\usepackage{float}
\usepackage{fancyhdr}
\usepackage{enumitem}
\usepackage{graphicx}
\usepackage{cleveref}
\usepackage{tikz-cd}
\usepackage{mathrsfs}

\geometry{a4paper, left=0.5in, top=0.5in, right=0.5in, bottom=0.7in}
\onehalfspacing

\newtheorem{theorem}{Theorem}[section]
\newtheorem{lemma}{Lemma}[section]
\newtheorem{corollary}{Corollary}[theorem]
\theoremstyle{definition}
\newtheorem{definition}{Definition}[section]
\newtheorem{example}{Example}[section]
\theoremstyle{remark}
\newtheorem*{remark}{\textbf{Remark}}

\newcommand{\e}{\varepsilon}
\newcommand{\st}{\text{ such that }}
\newcommand{\B}{\mathscr{B}}
\DeclareMathOperator{\Span}{span}
\DeclareMathOperator{\nullity}{nullity}
\DeclareMathOperator{\rank}{rank}

\hyphenpenalty=10000
\exhyphenpenalty=10000

\title{\textbf{\Large Structural Identity in Linear Spaces \\ \large A Rigorous Approach to Invertibility and Coordinate-Free Geometry}}
\author{\textbf{Tushar Kumar Aich}}
\date{\today}

\begin{document}
\maketitle

\begin{abstract}
    This paper presents a rigorous study of Invertibility and Isomorphism in linear algebra developing a coordinate-free perspective letting us understand that properties of a linear space remains independent of the chosen basis to describe it. We begin by examining the conditions under which a linear transformation $T: V \to W$ between vector spaces is invertible. Building upon this foundation, we introduce the concept of an isomorphism—a bijective linear map that acts as a structural bridge between vector spaces.
\end{abstract}

\section{Introduction}
Linear algebra studies vector spaces and the transformations that preserve their structure. But what happens when a linear map $T: V \to W$ can be perfectly reversed. Such an invertible map—which must be both one-to-one and onto—is known as an isomorphism. Isomorphism tells us that any two finite-dimensional vector spaces are algebraically identical if and only if they share the same dimension.

\section{Invertibility}

\begin{definition}
    Let $V$ and $W$ be vector spaces, and let $T: V \to W$ be linear. A function $U: W \to V$ is said to be an inverse of $T$ if $TU = I_W$ and $UT = I_V$. If $T$ has an inverse, then $T$ is said to be invertible. If $T$ is invertible, then the inverse of $T$ is unique and is denoted by $T^{-1}$.
\end{definition}
The following statements hold true for invertible functions $T$ and $U$:
\begin{enumerate}
    \item[(1)] $(TU)^{-1} = U^{-1}T^{-1}$
    \item[(2)] $(T^{-1})^{-1} = T$; i.e., $T^{-1}$ is invertible.
\end{enumerate}

We often use that a function is invertible if and only if it is both one-to-one and onto.

\begin{enumerate}
    \item[(3)] Let $T: V \to W$ be linear, where $V$ and $W$ are finite-dimensional spaces of equal dimension. Then $T$ is invertible if and only if $\rank(T) = \dim(V)$.
\end{enumerate}

\begin{theorem}
    Let $V$ and $W$ be vector spaces, and let $T: V \to W$ be linear and invertible. Then $T^{-1}: W \to V$ is linear.
\end{theorem}

\begin{proof}
    Let $y_1, y_2 \in W$ and $c \in F$. Since inverse of $T$ is defined then $T$ is invertible and thus $T$ is both one-to-one and onto. Therefore, there exists unique vectors $x_1$ and $x_2$ such that $T(x_1) = y_1$ and $T(x_2) = y_2$. Thus, $x_1 = T^{-1}(y_1)$ and $x_2 = T^{-1}(y_2)$. Thus,

    \begin{align*}
        T^{-1}(cy_1 + y_2) &= T^{-1}[cT(x_1) + T(x_2)] \\
        &= T^{-1}[T(cx_1 + x_2)] \\
        &= cx_1 + x_2 \\
        &= cT^{-1}(y_1) + T^{-1}(y_2)
    \end{align*}
\end{proof}

\begin{definition}
    Let $A$ be an $n \times n$ matrix. Then $A$ is invertible if there exists an $n \times n$ matrix $B$ such that $AB = BA = I$.
\end{definition}

 If $A$ is invertible then the matrix $B$ such that $AB = BA = I$ is unique. The matrix $B$ is called inverse of $A$ and is denoted by $A^{-1}$.

 \begin{lemma}
     Let $T$ be an invertible linear transformation from $V$ to $W$. Then $V$ is finite-dimensional if and only if $W$ is finite-dimensional. In this case, $\dim(V) = \dim(W)$.
 \end{lemma}

 \begin{proof}
     Since, $T^{-1}$ is defined from $W$ to $V$. Then $T$ is defined from $V$ to $W$. Hence,
     \[
     T^{-1}: W \to V \quad \text{and} \quad T: V \to W
     \]

     Also $T$ is one-to-one and onto. Thus $R(T) = W$. Now let $V$ be finite-dimensional and let $\B = \{x_1, x_2, \dots, x_n\}$ be a basis for $V$. By Theorem-2.2, $\text{span}(T(\B)) = R(T) = W$. By Theorem-1.6, $W$ is finite dimensional.

     Now we have $T^{-1}: W \to V$. Since $T^{-1}$ is always invertible. Thus $T^{-1}$ is also one-to-one and onto. Then $R(T^{-1}) = V$. If $W$ is finite dimensional and $\mathscr{C} = \{w_1, w_2, \dots, w_n\}$ is a finite basis for $W$. Then by Theorem-2.2, $\text{span}(T^{-1}(\mathscr{C})) = R(T^{-1}) = V$. Since $V$ is generated by a finite set $T^{-1}(\mathscr{C})$, then by Theorem-1.6, $V$ has a finite basis or we can say $V$ is finite-dimensional.

     Now since $T$ is one-to-one and onto. Thus,
     \[\dim(N(T)) = 0 \quad \text{and} \quad \dim(R(T)) = \dim(W).\]

     Thus, by rank-nullity theorem,
     \[\dim(V) = \dim(W)\]
 \end{proof}

 \begin{theorem}
     Let $V$ and $W$ be finite-dimensional vector spaces with ordered bases $\mathscr{B}$ and $\mathscr{C}$, respectively. Let $T: V \to W$ be linear. Then $T$ is invertible if and only if $[T]_{\mathscr{B}}^{\mathscr{C}}$ is invertible. Furthermore,
     $$[T^{-1}]_{\mathscr{C}}^{\mathscr{B}} = ([T]_{\mathscr{B}}^{\mathscr{C}})^{-1}$$
 \end{theorem}

 \begin{proof}
     Suppose $T$ is invertible. By the lemma 2.1, we have $\dim(V) = \dim(W)$. Let $n = \dim(V) = \dim(W)$. So $[T]_{\mathscr{B}}^{\mathscr{C}}$ is an $n \times n$ matrix. Now, $T^{-1}: W \to V$ satisfies $TT^{-1} = I_W$ and $T^{-1}T = I_V$. We know $I_n = [I_V]_{\mathscr{B}}$, thus we have:
     $$I_n = [I_V]_{\mathscr{B}} = [T^{-1}T]_{\mathscr{B}} = [T^{-1}]_{\mathscr{C}}^{\mathscr{B}} [T]_{\mathscr{B}}^{\mathscr{C}}$$
     Similarly, $[T]_{\mathscr{B}}^{\mathscr{C}} [T^{-1}]_{\mathscr{C}}^{\mathscr{B}} = I_n$. So $[T]_{\mathscr{B}}^{\mathscr{C}}$ is invertible and
     $$([T]_{\mathscr{B}}^{\mathscr{C}})^{-1} = [T^{-1}]_{\mathscr{C}}^{\mathscr{B}}$$

     Now suppose $A = [T]_{\mathscr{B}}^{\mathscr{C}}$ is invertible. Then there exists an $n \times n$ matrix $B$ such that $AB = BA = I_n$. There exists $U \in \mathscr{L}(W, V)$ such that
     $$U(w_j) = \sum_{i=1}^{n} B_{ij} v_i \quad \text{for } 1 \le j \le n$$
     where $\mathscr{C} = \{w_1, w_2, \dots, w_n\}$ and $\mathscr{B} = \{v_1, v_2, \dots, v_n\}$. It shows that $[U]_{\mathscr{C}}^{\mathscr{B}} = B$. Again,
     $$[UT]_{\mathscr{B}}^{\mathscr{B}} = [U]_{\mathscr{C}}^{\mathscr{B}} [T]_{\mathscr{B}}^{\mathscr{C}} = BA = I_n = [I_V]_{\mathscr{B}}$$
     Thus, $UT = I_V$ and similarly, $TU = I_W$. Thus, $U$ is inverse of $T$ or $U = T^{-1}$. Hence, $[T]_{\mathscr{B}}^{\mathscr{C}}$ is invertible and,
     $$([T]_{\mathscr{B}}^{\mathscr{C}})^{-1} = [U]_{\mathscr{C}}^{\mathscr{B}} = [T^{-1}]_{\mathscr{C}}^{\mathscr{B}}$$
 \end{proof}

 \begin{corollary}
     Let $V$ be a finite-dimensional vector space with ordered basis $\mathscr{B}$, and let $T: V \to V$ be linear. Then $T$ is invertible if and only if $[T]_{\mathscr{B}}$ is invertible. Furthermore,
     $$[T]_{\mathscr{B}}^{-1} = [T^{-1}]_{\mathscr{B}}$$
 \end{corollary}

 \begin{proof}
     From Theorem 2.2, any transformation $T: V \to W$ with bases $\mathscr{B}$ and $\mathscr{C}$, $T$ is invertible if and only if $[T]_{\mathscr{B}}^{\mathscr{C}}$ is invertible. By setting $W = V$ and $\mathscr{B} = \mathscr{C}$, it follows immediately that $T: V \to V$ is invertible if and only if $[T]_{\mathscr{B}}$ is invertible. Also we had,
$$[T^{-1}]_{\mathscr{C}}^{\mathscr{B}} = ([T]_{\mathscr{B}}^{\mathscr{C}})^{-1}$$
Again by setting $W = V$ and $\mathscr{B} = \mathscr{C}$, we get,
$$[T^{-1}]_{\mathscr{B}} = ([T]_{\mathscr{B}})^{-1}$$
\end{proof}

\begin{corollary}
    Let $A$ be an $n \times n$ matrix. Then $A$ is invertible if and only if $L_A$ is invertible. Furthermore, $(L_A)^{-1} = L_{A^{-1}}$.
    
\end{corollary}

\begin{proof}
    By Corollary 2.2.1, $L_A$ is invertible if and only if $[L_A]_{\mathscr{B}}$ is invertible. We know if $A$ is an $n \times n$ matrix and $\beta$ is a standard ordered basis for $F^n$ and $L_A:F^n \to F^n$ is linear then $[L_A]_\beta = A$, we can say that $[L_A]_{\mathscr{B}} = A$. So now $L_A$ is invertible if and only if $A$ is invertible.

Now to prove $(L_A)^{-1} = L_{A^{-1}}$, we must satisfy that
$$L_A \cdot L_{A^{-1}} = I_{F^n} \quad \text{and} \quad L_{A^{-1}} \cdot L_A = I_{F^n}$$
we know for another $n \times n$ matrix $B$, $L_{AB} = L_A . L_B$, then
$$L_A \cdot L_{A^{-1}} = L_{AA^{-1}} \quad \text{and} \quad L_{A^{-1}} \cdot L_A = L_{A^{-1}A}$$
By definition of inverse of a matrix,
$$L_{AA^{-1}} = L_{I_n} \quad \text{and} \quad L_{A^{-1}A} = L_{I_n}$$
We know $L_{I_n} = I_{F^n}$,
Thus, $L_A \cdot L_{A^{-1}} = I_{F^n}$ and $L_{A^{-1}} \cdot L_A = I_{F^n}$.
\end{proof}

\begin{theorem}
    Let $A$ and $B$ be $n \times n$ invertible matrices. Prove that $AB$ is invertible and $(AB)^{-1} = B^{-1}A^{-1}$.
\end{theorem}

\begin{proof}
    Let $V, W,$ and $Z$ be $n$-dimensional vector spaces. Let $T: V \to W$ and $U: W \to Z$ be linear, and let $\alpha, \beta,$ and $\gamma$ be the ordered bases of $V, W,$ and $Z$ respectively, such that $A = [U]_\beta^\gamma$ and $B = [T]_\alpha^\beta$. Thus $AB = [UT]_\alpha^\gamma$.

Since $A$ is invertible, then by Theorem 2.2, $U$ is invertible. Thus, $U$ is one-to-one and onto; similarly, $T$ is one-to-one and onto. We know that if $U$ and $T$ are both one-to-one and onto, then $UT$ is also one-to-one and onto. Thus, $UT$ is invertible and, by Theorem 2.2, $[UT]_\alpha^\gamma$ or $AB$ is invertible. Thus,
\[ (AB)^{-1} = ([UT]_\alpha^\gamma)^{-1} = [(UT)^{-1}]_\gamma^\alpha \]
We are given, for invertibility, $(UT)^{-1} = T^{-1}U^{-1}$. Thus,
\[ (AB)^{-1} = [(UT)^{-1}]_\gamma^\alpha = [T^{-1}U^{-1}]_\gamma^\alpha = [T^{-1}]_\beta^\alpha [U^{-1}]_\gamma^\beta = B^{-1}A^{-1} \]
\end{proof}

\section{Isomorphism}
\begin{definition}
    Let $V$ and $W$ be vector spaces. We can say $V$ \textit{is isomorphic to} $W$ if there exists a linear transformation $T:V\to W$ that is invertible. Such a linear transformation is called an isomorphism from $V$ onto $W$.
\end{definition}

\begin{theorem}
    Let $V$ and $W$ be finite-dimensional vector spaces over the same field $F$. Then $V$ \textit{is isomorphic to} $W$ if and only if $\dim(V) = \dim(W)$.
\end{theorem}

\begin{proof}
    Suppose $V$ is isomorphic to $W$, then $T: V \to W$ is an isomorphism and thus, by the preceding Theorem 2.2, $\dim(V) = \dim(W)$.

    Suppose $\dim(V) = \dim(W)$, let $\beta = \{v_1, v_2, \dots, v_n\}$ and $\mathcal{C} = \{w_1, w_2, \dots, w_n\}$ be bases for $V$ and $W$ respectively. Now there exists $T: V \to W$ such that $T$ is linear and $T(v_i) = w_i$ for $i = 1, 2, \dots, n$. We know if $T:V \to W$ is linear then, $R(T) = \text{span}(\mathcal{C}) = W$. So $T$ is onto. Now by rank-nullity theorem, we have
    \begin{align*}
        \dim(V) &= \dim(R(T)) + \dim(N(T)) \\
        \implies \dim(W) &= \dim(W) + \dim(N(T)) \\
        \implies \dim(N(T)) &= 0
    \end{align*}
    Thus $N(T) = \{0\} $ and therefore $T$ is one-to-one. Hence $T$ is an isomorphism and hence $V$ \textit{is isomorphic to} $W$.
\end{proof}

This theorem proves that equal dimension is the key to structural equivalence. Because the dimensions are equal, both spaces have the same number of basis vectors, which allows us to create a direct correspondence by mapping the basis of one space to the basis of the other.

Since the dimensions match, there is no "excess" space in one that would cause information loss, and no "lack" of space that would prevent us from reaching every vector. This guarantees that our transformation is both one-to-one and onto—essentially a perfect, reversible bridge. Now, whenever we want to know if we can move between two vector spaces without losing any information, we just have to check if their dimensions are equal.

\begin{corollary}
    Let $V$ be a vector space over $F$. Then $V$ \textit{is isomorphic to} $F^n$ if and only if $\dim(V) = n$.
\end{corollary}

\begin{proof}
    From Theorem 3.1, we have $V$ \textit{is isomorphic to} $W$ if and only if $\dim(V) = \dim(W)$. Let $W = F^n$, we get $V$ \textit{is isomorphic to} $F^n$ if and only if $\dim(V) = \dim(F^n)$. We know $\dim(F^n) = n$. Thus we get $V$ \textit{is isomorphic to} $F^n$ if and only if $\dim(V) = n$.
\end{proof}

We’ve realized that every $n$-dimensional vector space $V$ is structurally equivalent to $F^n$. This means we can represent any abstract vector in $V$ as a column matrix in $F^n$ relative to a chosen basis. This is the fundamental link that connects abstract linear maps to matrix algebra which allows us to perform computations in $F^n$ while knowing that the structural properties are perfectly preserved back in $V$.

\begin{theorem}
    Let $V$ and $W$ be finite-dimensional vector spaces over $F$ of dimensions $n$ and $m$, respectively, and let $\B$ and $\mathcal{C}$ be ordered bases for $V$ and $W$, respectively. Then $\Phi: \mathscr{L}(V, W) \to M_{m \times n}(F)$, defined by $\Phi(T) = [T]_\B^\mathscr{C}$ for $T \in \mathscr{L}(V, W)$, is an isomorphism.
\end{theorem}

\begin{proof}
    Let $T_1, T_2 \in \mathscr{L}(V, W)$ and $a \in F$. Then, for linearity, we need to show $\Phi(aT_1 + T_2) = a\Phi(T_1) + \Phi(T_2)$. Now,
\[ \Phi(aT_1 + T_2) = [aT_1 + T_2]_\B^\mathcal{C} \]
\[ [aT_1 + T_2]_\B^\mathcal{C} = a[T_1]_\B^\mathcal{C} + [T_2]_\B^\mathcal{C} = a\Phi(T_1) + \Phi(T_2) \]
Hence, $\Phi$ is linear. Now to prove isomorphism, we need to show that $\Phi$ is one-to-one and onto. Let $\Phi(T) = 0_{M_{m \times n}(F)}$, so $[T]_\B^\mathcal{C} = 0_{M_{m \times n}(F)}$. Here, the RHS is a matrix of order $m \times n$ with all its entries $0$. Thus, the LHS can be equal to the RHS only if all the entries of $[T]_\B^\mathcal{C}$ are $0$; in other words, $T = 0_{\mathscr{L}(V, W)}$. Thus, $N(\Phi) = \{0_{\mathscr{L}(V, W)}\}$. Hence, $\Phi$ is one-to-one.

Now, let $A \in M_{m \times n}(F)$ with entries $A_{ij}$. Now there exists a unique transformation $T: V \to W$ such that $T(v_j) = \sum_{i=1}^{n} A_{ij} w_i$ for $1 \leq j \leq n$. By definition, $\Phi(T) = [T]_\B^\mathcal{C} = A$. Thus, $\Phi$ is onto. By definition, $\Phi$ is invertible, and hence $\Phi$ is an isomorphism.
\end{proof}

\begin{corollary}
    Let $V$ and $W$ be finite-dimensional vector spaces of dimensions $n$ and $m$ respectively. Then $\mathscr{L}(V, W)$ is finite-dimensional of dimension $mn$.
\end{corollary}

\begin{proof}
    By Theorem 3.2, $\Phi: \mathscr{L}(V, W) \to M_{m \times n}(F)$ is an isomorphism. By Theorem 3.1, $\dim(\mathscr{L}(V, W)) = \dim(M_{m \times n}(F))$. We know $\dim(M_{m \times n}(F)) = mn$. Thus, $\dim(\mathscr{L}(V, W)) = mn$.
\end{proof}

So, as previously we showed that every vector in a finite-dimensional vector space $V$ can be represented as a column matrix in $F^n$, similarly here we can say that any transformation $T$ between the two finite-dimensional vector spaces $V$ and $W$ can be represented as a matrix in $M_{m \times n}(F)$. This is the complete bridge between the abstract world of maps and the computational world of matrices. It confirms that the entire machinery of linear algebra both vectors and the maps between them can be fully expressed through matrix algebra without sacrificing any information.

\begin{theorem}
    Let $B$ be an $n \times n$ invertible matrix. Define $\Phi: M_{n \times n}(F) \to M_{n \times n}(F)$ by $\Phi(A) = B^{-1}AB$. Then $\Phi$ is an isomorphism.
\end{theorem}

\begin{proof}
    For $A_1, A_2 \in M_{n \times n}(F)$ and $c \in F$, using the distributive and associative properties of matrix multiplication, we have:
    \[ \Phi(cA_1 + A_2) = B^{-1}(cA_1 + A_2)B = c(B^{-1}A_1B) + (B^{-1}A_2B) = c\Phi(A_1) + \Phi(A_2) \]
    Thus, $\Phi$ is linear.

    Now let $A \in N(\Phi)$, then $\Phi(A) = 0 \st 0 = 0_{M_{n \times n}(F)}$. Multiplying by $B$ on the left and $B^{-1}$ on the right, we get
    \[ B(B^{-1}AB)B^{-1} = B(0)B^{-1} \implies A = 0 \]
    Thus $N(\Phi) = \{0_{M_{n \times n}(F)}\}$. Hence $\Phi$ is one-to-one.

    Again by rank-nullity theorem,
    \[\dim(M_{n \times n}(F)) = \dim(R(\Phi)) + \dim(N(\Phi)) \implies \dim(M_{n \times n}(F)) = \dim(R(\Phi))\]
    Since, $R(\Phi)$ is a subspace of $M_{n \times n}(F)$ then $R(\Phi) = M_{n \times n}(F)$. Therefore, $\Phi$ is onto. Hence $\Phi$ is an isomorphism.
\end{proof}

\begin{definition}
    Let $\B$ be an ordered basis for an $n$-dimensional vector space $V$ over the field $F$. The standard representation of $V$ with respect to $\B$ is the transformation $\phi_\B:V \to F^n$ defined by $\phi_\B(x) = [x]_\B$ for all $x \in V$.
\end{definition}

\begin{theorem}
    For any finite-dimensional vector space $V$ with ordered basis $\B$, $\phi_\B$ is an isomorphism.
\end{theorem}

\begin{proof}
    Let $x, y \in V$ and $c \in F$. Since $\B$ is a basis for $V$, then for $\B = \{v_1, v_2, \dots, v_n\}$, there exist scalars $a_i, b_i \in F$ for $i = 1, 2, \dots, n$ such that:
\[ x = \sum_{i=1}^{n} a_i v_i \quad \text{and} \quad y = \sum_{i=1}^{n} b_i v_i \]
Thus, $cx + y = \displaystyle\sum_{i=1}^{n} (ca_i + b_i) v_i$. Therefore,
\[ \phi_\B(cx + y) = \begin{bmatrix} ca_1 + b_1 \\ ca_2 + b_2 \\ \vdots \\ ca_n + b_n \end{bmatrix} = c \begin{bmatrix} a_1 \\ a_2 \\ \vdots \\ a_n \end{bmatrix} + \begin{bmatrix} b_1 \\ b_2 \\ \vdots \\ b_n \end{bmatrix} = c \phi_\B(x) + \phi_\B(y) \]
Hence, $\phi_\B$ is linear. Now, let $x \in N(\phi_\B)$:
\[ \phi_\B(x) = 0_{F^n} \implies [x]_\B = 0_{F^n} \]
Thus, the LHS is a vector consisting only of zeroes. Therefore, $x = 0_V$. Hence, $N(\phi_\B) = \{0_V\}$ and $\phi_\B$ is one-to-one. This implies $\dim(N(\phi_\B)) = 0$. Now, by the Rank-Nullity Theorem:
\[ \dim(V) = \dim(R(\phi_\B)) + \dim(N(\phi_\B)) \]
We have $\dim(V) = n$. Thus,
\[ \dim(R(\phi_\B)) = n = \dim(F^n) \]
But since $R(\phi_\B)$ is a subspace of $F^n$, then $R(\phi_\B) = F^n$. Therefore, $\phi_\B$ is onto, and hence $\phi_\B$ is an isomorphism.
\end{proof}

This confirms our previous intuition that coordinate-maps are not just a computational tool but an isomorphism that preserves the structure of the vector space $V$ within the familiar environment if $F^n$.

\begin{figure}[h]
    \centering
    \begin{tikzcd}[row sep=huge, column sep=2cm, labels={font=\large}]
        V \arrow[r, "T"] \arrow[d, "\phi_\mathscr{B}"]
        & W \arrow[d, "\phi_\mathscr{C}"] \\
        F^n \arrow[r, "L_A"]
        & F^m
    \end{tikzcd}
\end{figure}

Here there are two composites of linear transformations from $V$ to $F^m$.
\begin{enumerate}
    \item Map $V$ to $F^n$ with $\phi_\B$ and then follow this transformation with $L_A$, \textit{i.e.}, $L_A\phi_\B$.

    \item Map $V$ to $W$ with $T$ and then follow this with $\phi_\mathscr{C}$ to obtain $\phi_\mathscr{C}T$.
   
\end{enumerate}
Thus, $L_A\phi_\B = \phi_\mathscr{C}T$.

\begin{lemma}
    Let $V$ and $W$ be finite-dimensional vector spaces and $T:V \to W$ be an isomorphism. Let $V_0$ be a subspace of $V$ then $\dim(V_0) = \dim(T(V_0))$.
\end{lemma}

\begin{proof}
    Since $V$ and $W$ are finite-dimensional, $V_0$ and $T(V_0)$ are also finite-dimensional. Let $\B = \{v_1, v_2, \dots, v_k\}$ be a basis for $V_0$. Then $T(\B) = \{T(v_1), T(v_2), \dots, T(v_k)\}$.

\textbf{Proving $T(\B)$ spans $T(V_0)$:}
Let $y \in T(V_0)$, then there exists $x \in V_0$ such that $T(x) = y$. Since $\B$ is a basis for $V_0$, we can write $x$ as:
\[ x = \sum_{i=1}^{k} c_i v_i \quad \text{for } c_i \in F \]
Applying $T$ to both sides:
\[ T(x) = T\left(\sum_{i=1}^{k} c_i v_i\right) = \sum_{i=1}^{k} c_i T(v_i) \]
\[ \implies y = \sum_{i=1}^{k} c_i T(v_i) \]
Thus, any $y \in T(V_0)$ can be expressed as a linear combination of $T(\B)$. Hence, $\text{span}(T(\B)) = T(V_0)$.

\textbf{Proving $T(\B)$ is linearly independent:}
Let a linear combination of vectors in $T(\B)$ be equal to $0_W$:
\[ \sum_{i=1}^{k} c_i T(v_i) = 0_W \]
By the linearity of $T$:
\[ T\left(\sum_{i=1}^{k} c_i v_i\right) = 0_W \]
Therefore, $\displaystyle\sum_{i=1}^{k} c_i v_i \in N(T)$. Since $T$ is one-to-one, $N(T) = \{0_V\}$. Thus,
\[ \sum_{i=1}^{k} c_i v_i = 0_V \]
Since $\B$ is a basis for $V_0$, the scalars $c_i = 0$ for all $i = 1, 2, \dots, k$. Hence, $T(\B)$ is linearly independent.

Since $T(\B)$ is a basis for $T(V_0)$ and contains $k$ vectors, $\dim(T(V_0)) = k$. Also, $\dim(V_0) = k$. Therefore:
\[ \dim(V_0) = \dim(T(V_0)) \]
\end{proof}

\begin{theorem}
    Let $T: V \to W$ be a linear transformation from an $n$-dimensional vector space $V$ to an $m$-dimensional vector space $W$. Let $\beta$ and $\gamma$ be ordered bases for $V$ and $W$ respectively. Prove that $\text{rank}(T) = \text{rank}(L_A)$ and that $\text{nullity}(T) = \text{nullity}(L_A)$, where $A = [T]_\beta^\gamma$.
\end{theorem}

\begin{proof}
    Let us define transformations $\phi_\beta: V \to F^n$, $L_A: F^n \to F^m$, and $\phi_\gamma: W \to F^m$ such that $L_A \phi_\beta = \phi_\gamma T$. From Theorem 3.4, we know $\phi_\beta$ and $\phi_\gamma$ are isomorphisms. Thus, $\dim(V) = \dim(F^n)$ and $\dim(W) = \dim(F^m)$.

Now, since $T$ and $L_A$ are linear transformations, by the Rank-Nullity Theorem:
\[ \dim(V) = \text{rank}(T) + \text{nullity}(T) \quad \text{and} \quad \dim(F^n) = \text{rank}(L_A) + \text{nullity}(L_A) \]
Since $\dim(V) = \dim(F^n)$, we have:
\[ \text{rank}(T) + \text{nullity}(T) = \text{rank}(L_A) + \text{nullity}(L_A) \]
\[ \implies \text{rank}(T) - \text{rank}(L_A) = \text{nullity}(L_A) - \text{nullity}(T) \tag{1} \]

Again, since $R(T)$ is a subspace of $W$ and $\phi_\gamma$ is an isomorphism, then by lemma 3.1 $\dim(R(T)) = \dim(\phi_\gamma(R(T)))$, or $\text{rank}(T) = \dim(\phi_\gamma(R(T)))$. We know $R(T)$ is the mapping of all vectors in $V$ through $T(V)$. Thus $R(T) = T(V)$. Therefore, $\phi_\gamma(R(T)) = \phi_\gamma(T(V))$. Since $\phi_\gamma T = L_A \phi_\beta$, we have $\phi_\gamma(R(T)) = L_A(\phi_\beta(V))$. Since $\phi_\beta$ is an isomorphism, $\phi_\beta(V) = F^n$. So, $\phi_\gamma(R(T)) = L_A(F^n)$. 

We know $R(L_A)$ is the mapping of all vectors in $F^n$ through $L_A$, thus $R(L_A) = L_A(F^n)$. So, $\phi_\gamma(R(T)) = R(L_A)$, then $\dim(\phi_\gamma(R(T))) = \dim(R(L_A))$, or $\text{rank}(L_A) = \text{rank}(T)$. Thus, $\text{rank}(T) - \text{rank}(L_A) = 0$.

Substituting this into equation (1), we get $\text{nullity}(L_A) - \text{nullity}(T) = 0$, which implies $\text{nullity}(T) = \text{nullity}(L_A)$.
\end{proof}

\begin{theorem}[\textbf{The First Isomorphism Theorem}]
    Let $T: V \to Z$ be a linear transformation for a vector space $V$ onto a vector space $Z$. Define the mapping $\overline{T}:V/N(T) \to Z$ by $\overline{T}(v+N(T))=T(v)$ for any coset $v+N(T) \in V/N(T)$. Then
    \begin{enumerate}
        \item $\overline{T}$ is well defined, \textit{i.e.}, if $v+N(T) = v'+N(T)$ then $T(v)=T(v')$.
        \item $\overline{T}$ is linear.
        \item $\overline{T}$ is an isomorphism.
        \item $T = \overline{T}\eta$
    \end{enumerate}
    \begin{figure}[h]
        \centering
        \begin{tikzcd}[row sep=huge, column sep=2cm, labels={font=\large}]
            V \arrow[r, "T"] \arrow[d, "\overline{T}"]
            & Z \\
            V/N(T) \arrow[ur, "\eta"]
        \end{tikzcd}
    \end{figure}
\end{theorem}

\begin{proof}
    \begin{enumerate}
        \item We know if $v+N(T) = v'+N(T)$ for any coset in $V/N(T)$ then $v - v' \in N(T)$. By definition of $N(T)$, $T(v-v') = 0_Z$. Since $T$ is linear, $T(v-v') = T(v) - T(v') = 0_Z$ which implies $T(v) = 0_Z + T(v')$. Since $Z$ is a subspace, $T(v) = 0_Z + T(v') = T(v')$.

        \item Let $v_1 + N(T), v_2 + N(T) \in V/N(T)$ and $c \in F$. We need to prove:
        \[ \overline{T}[c(v_1 + N(T)) + (v_2 + N(T))] = c\overline{T}(v_1 + N(T)) + \overline{T}(v_2 + N(T)) \]

        Using the properties of cosets, we know that $(v_1 + W) + (v_2 + W) = (v_1 + v_2) + W$ and $c(v_1 + W) = cv_1 + W$. Thus:
        \[ c(v_1 + N(T)) + (v_2 + N(T)) = (cv_1 + v_2) + N(T) \]

        Therefore:
        \[ \overline{T}[(cv_1 + v_2) + N(T)] = T(cv_1 + v_2) \]
        Since $T$ is linear:
        \[ T(cv_1 + v_2) = cT(v_1) + T(v_2) = c\overline{T}(v_1 + N(T)) + \overline{T}(v_2 + N(T)) \]
        Hence, $\overline{T}$ is linear.
        \item To prove $\overline{T}$ is one-to-one, we need to show that $N(\overline{T}) = \{0_{V/N(T)}\} = \{N(T)\}$. Let an arbitrary coset $v + N(T) \in N(\overline{T})$ for any $v \in V$. Then:
        \[ \overline{T}(v + N(T)) = 0_Z \]
        \[ \implies T(v) = 0_Z \]
        Since $T(v) = 0_Z$, then by definition $v \in N(T)$. Therefore, the coset $v + N(T)$ is equal to the zero coset $N(T)$. Thus, $N(\overline{T}) = \{N(T)\}$, which proves that $\overline{T}$ is one-to-one.

        Since $R(\overline{T})$ is a subspace of $Z$, it follows that $R(\overline{T}) \subseteq Z$. Now, let $z \in Z$. Since $T$ is onto, there exists some $v \in V$ such that $T(v) = z$. Since $T(v) = \overline{T}(v + N(T))$, we have:
        \[ \overline{T}(v + N(T)) = z \]
        Therefore, $z \in R(\overline{T})$, which implies $Z \subseteq R(\overline{T})$. Since $R(\overline{T}) \subseteq Z$ and $Z \subseteq R(\overline{T})$, we have $R(\overline{T}) = Z$, and hence $\overline{T}$ is onto.

        Since $\overline{T}$ is both one-to-one and onto, it is an isomorphism.

        \item We define $\eta:V \to V/N(T)$ by $\eta(v) = v + N(T) \text{ for } v \in V$. Thus $\overline{T}\eta:V \to Z$. Let $v \in V$ then,
        \[\overline{T}\eta(v) = \overline{T}(\eta(v)) = \overline{T}(v + N(T)) = T(v)\]
        Thus, $\overline{T}\eta = T$.
    \end{enumerate}
\end{proof}

\section{Conclusion}
In these notes, we’ve bridged the gap between abstract vector spaces and their computational matrices. We proved that for any finite-dimensional space, choosing a basis creates an isomorphism to $F^n$, effectively turning geometry into matrix arithmetic. Furthermore, the First Isomorphism Theorem clarified that the structure of any linear map T is essentially an isomorphism acting on the quotient space V/N(T).

Ultimately, these results show that matrices aren't just for calculation—they are the faithful language of linear structure. Whether we are changing bases or analyzing transformations, we can trust that these isomorphisms preserve the underlying geometric truth.

\begin{thebibliography}{9}
    \bibitem{fis_linear_algebra}
    Friedberg, S. H., Insel, A. J., \& Spence, L. E. (2018). \textit{Linear Algebra} (5th ed.). Pearson.
\end{thebibliography}
\end{document}
